% !TEX program = xelatex
\documentclass[11pt,a4paper]{article}
\usepackage[a4paper,margin=18mm,headheight=15pt]{geometry}
\usepackage{fontspec}\usepackage{polyglossia}\setdefaultlanguage{russian}\setotherlanguage{english}
\setmainfont{Arial}\setsansfont{Arial}\setmonofont{Arial}
\usepackage{microtype,enumitem,array,booktabs,tabularx,xcolor,hyperref,fancyhdr,titlesec,tcolorbox,amssymb,lastpage}
\definecolor{MaxPurple}{HTML}{641CF3}\definecolor{MaxBlue}{HTML}{0A84FF}\definecolor{Ink}{HTML}{16181D}\definecolor{Muted}{HTML}{5E6673}\definecolor{Panel}{HTML}{F3F6FA}\definecolor{Alert}{HTML}{D92D20}
\hypersetup{colorlinks=true,linkcolor=MaxPurple,urlcolor=MaxBlue,pdftitle={Спринт 2 - вертикальный MVP},pdfauthor={Команда QR-посещаемости в MAX}}
\pagestyle{fancy}\fancyhf{}\fancyhead[L]{\small\color{Muted}QR-посещаемость в MAX}\fancyhead[R]{\small\color{Muted}Спринт 2 \textbullet{} 24-25 сентября}\fancyfoot[L]{\small\color{Muted}План команды из 4 человек}\fancyfoot[R]{\small\color{Muted}\thepage\ / \pageref{LastPage}}\renewcommand{\headrulewidth}{0.4pt}
\titleformat{\section}{\Large\bfseries\color{Ink}}{\colorbox{MaxPurple}{\textcolor{white}{\strut\thesection}}}{0.65em}{}\titleformat{\subsection}{\large\bfseries\color{Ink}}{\thesubsection}{0.55em}{}\titlespacing*{\section}{0pt}{1.9ex}{0.9ex}\titlespacing*{\subsection}{0pt}{1.3ex}{0.5ex}
\newlist{tasklist}{itemize}{1}\setlist[tasklist]{label=$\square$,leftmargin=1.6em,itemsep=0.22em,topsep=0.3em}\setlist[itemize]{leftmargin=1.55em,itemsep=0.2em,topsep=0.3em}
\setlength{\parindent}{0pt}
\newcommand{\code}[1]{\texttt{#1}}\newcommand{\owner}[1]{\textcolor{MaxPurple}{\textbf{Владелец: #1}}}
\newtcolorbox{goalbox}[1][]{colback=Panel,colframe=MaxPurple,boxrule=0.9pt,arc=2mm,left=3mm,right=3mm,top=2.5mm,bottom=2.5mm,#1}\newtcolorbox{alertbox}[1][]{colback=white,colframe=Alert,boxrule=0.9pt,arc=2mm,left=3mm,right=3mm,top=2.5mm,bottom=2.5mm,#1}
\begin{document}
\begin{titlepage}\pagecolor{MaxPurple}\color{white}\vspace*{17mm}{\fontsize{15}{18}\selectfont ПРОЕКТ \textbullet{} QR-ПОСЕЩАЕМОСТЬ В MAX}\par\vspace{19mm}{\fontsize{36}{42}\selectfont\bfseries Спринт 2\\Вертикальный\\рабочий MVP}\par\vspace{9mm}{\Large 24-25 сентября 2026}\par\vfill
\begin{tcolorbox}[colback=white,colframe=white,arc=3mm,boxrule=0pt]\color{Ink}\textbf{Цель:} провести одного преподавателя и одного студента через полный сценарий без ручного вмешательства в базу: группа $\rightarrow$ сессия $\rightarrow$ QR $\rightarrow$ check-in $\rightarrow$ список $\rightarrow$ закрытие.\end{tcolorbox}\vspace{8mm}{\large Команда: Данил, Рома, Никита, Амир}\par\end{titlepage}\nopagecolor\color{Ink}

\section{Результат спринта}
\begin{goalbox}
К вечеру 25 сентября готов первый end-to-end MVP. Он может быть визуально неполированным, но обязан работать внутри MAX на реальных тестовых данных и через настоящие API. Демонстрация проходит без SQL-консоли, Postman и ручного изменения состояния.
\end{goalbox}
\begin{alertbox}\textbf{Правило спринта:} никакой работы над P2. Любой blocker основного сценария получает высший приоритет и владельца в день обнаружения.\end{alertbox}

\section{Распределение ответственности}
\begin{tabularx}{\textwidth}{>{\bfseries}p{25mm}p{34mm}X}\toprule
Участник & Основная роль & Результат к концу спринта\\\midrule
Данил & acceptance, QA, интеграция & проверяемые критерии и стабильный ежедневный demo build\\
Рома & product, UX, backend-support & группы, тексты интерфейса, UX-проверка и продуктовые слайды\\
Никита & frontend и MAX & весь путь преподавателя и студента в Mini App\\
Амир & backend и данные & полный API основного сценария и безопасный check-in\\\bottomrule
\end{tabularx}

\section{Задачи Данила}
\owner{Данил}
\begin{tasklist}
  \item Описать acceptance criteria основного сценария.
  \item Подготовить тестовые аккаунты.
  \item Подготовить тестовую группу.
  \item Создать список интеграционных проверок.
  \item Обновлять README по мере реализации.
  \item Ежедневно проводить общий интеграционный прогон.
  \item Не допускать добавления P2-функций.
  \item Фиксировать обнаруженные блокеры, владельцев и сроки исправления.
  \item Не принимать заглушки или ручное редактирование БД как готовую функцию.
\end{tasklist}

\section{Задачи Ромы}
\owner{Рома - product lead и backend-support}
\begin{tasklist}
  \item Подготовить тестовый состав группы.
  \item Реализовать или помочь реализовать endpoints групп.
  \item Реализовать привязку студентов к группе.
  \item Проверить тексты интерфейса.
  \item Подготовить сообщения об успешной отметке и ошибках.
  \item Проверить понятность сценария на пользователях.
  \item Зафиксировать результаты мини-UX-теста и выбрать blocker/P1-исправления.
  \item Начать материалы презентации: проблема, доказательства, As Is и To Be.
  \item Проверить соответствие сценария ценности и метрикам MVP.
\end{tasklist}

\section{Задачи Никиты}
\owner{Никита - frontend и MAX integration lead}
\begin{tasklist}
  \item Реализовать авторизацию через MAX.
  \item Реализовать экран выбора группы.
  \item Реализовать запуск attendance-сессии.
  \item Реализовать экран QR.
  \item Добавить таймер QR.
  \item Реализовать студенческий экран.
  \item Передавать \code{start\_param} backend-серверу.
  \item Реализовать подтверждение посещаемости.
  \item Реализовать живой список присутствующих.
  \item Реализовать завершение сессии.
  \item Реализовать итоговый экран.
  \item Проверить весь путь в мобильном и web-клиенте MAX.
\end{tasklist}

\section{Задачи Амира}
\owner{Амир - backend и security lead}
\begin{tasklist}
  \item Реализовать авторизацию MAX.
  \item Создавать и обновлять пользователя из валидного \code{initData}.
  \item Определять роль пользователя и проверять её на сервере.
  \item Реализовать группы и memberships.
  \item Реализовать создание attendance-сессии.
  \item Реализовать создание случайного QR-токена.
  \item Хранить только hash QR-токена.
  \item Добавить срок действия токена.
  \item Реализовать endpoint check-in.
  \item Проверять принадлежность студента к группе.
  \item Запретить повторную отметку.
  \item Возвращать список присутствующих.
  \item Реализовать закрытие сессии.
  \item Возвращать стабильные коды ошибок для frontend.
  \item Добавить интеграционный тест полного happy path.
\end{tasklist}

\section{Постоянные владельцы направлений}
\begin{tabularx}{\textwidth}{>{\bfseries}p{25mm}X X}\toprule
Участник & Владеет разделами & Поддерживает\\\midrule
Данил & управление, границы MVP, тестирование, документация, презентация, demo, сдача & все направления на приёмке\\
Рома & продуктовое исследование, пилот, масштабирование & MVP, UX, backend, frontend, презентация\\
Никита & UX/UI, MAX integration, frontend & архитектура, backend, документация\\
Амир & архитектура, backend, безопасность, Docker/DevOps & MVP, MAX integration, документация, сдача\\\bottomrule
\end{tabularx}

\section{Ежедневный ритм}
\begin{itemize}
  \item \textbf{10:00 - stand-up, 10 минут:} сделано, план, блокеры.
  \item \textbf{15:00 - быстрая интеграция:} frontend и backend проверяют совместимость.
  \item \textbf{20:00 - вечерняя демонстрация:} показывается только работающий результат.
  \item После демонстрации Данил обновляет приоритеты следующего дня.
\end{itemize}
\begin{alertbox}\textbf{Критический стык: Никита + Амир.} Frontend и backend интегрируются ежедневно; нельзя ждать окончания разработки отдельных веток.\end{alertbox}

\section{Основной сценарий приёмки}
\begin{enumerate}[leftmargin=1.7em,itemsep=0.25em]
  \item Преподаватель открывает Mini App в MAX и видит свои группы.
  \item Преподаватель открывает тестовую группу и запускает занятие.
  \item На экране появляется QR-код и активный таймер.
  \item Студент сканирует QR, переходит по deep link в MAX и подтверждает отметку.
  \item Backend валидирует пользователя, membership, сессию и QR-токен.
  \item Студент получает успешный результат, преподаватель видит его в списке.
  \item Повторная попытка не создаёт вторую запись.
  \item Преподаватель закрывает сессию и видит итог.
\end{enumerate}

\section{Точки интеграции}
\begin{tabularx}{\textwidth}{p{30mm}X>{\raggedright\arraybackslash}p{37mm}}\toprule
Когда & Совместная проверка & Участники\\\midrule
24 сен., 15:00 & auth, groups и memberships работают через UI & Рома + Никита + Амир\\
24 сен., 20:00 & создание сессии и показ реального QR & вся команда\\
25 сен., 15:00 & check-in, live-список, duplicate и close & Никита + Амир\\
25 сен., 20:00 & полный happy path на тестовых аккаунтах & вся команда\\\bottomrule
\end{tabularx}

\section{Критерии выхода}
\begin{tasklist}
  \item Полный happy path проходит внутри MAX без ручного вмешательства.
  \item Преподаватель может создать/выбрать группу и запустить сессию.
  \item QR содержит ограниченный по времени токен, а не идентификатор занятия в открытом виде.
  \item Участник группы отмечается ровно один раз.
  \item Неучастник группы и невалидный токен отклоняются сервером.
  \item Список присутствующих обновляется и сохраняется после закрытия сессии.
  \item Все API-вызовы выполняются frontend-приложением, данные хранятся в PostgreSQL.
  \item README позволяет другому участнику повторить демонстрацию.
\end{tasklist}

\section{Рабочий журнал}
\begin{tabularx}{\textwidth}{p{25mm}X X}\toprule Дата & Сделано / доказательство & Блокер / решение\\\midrule
24 сентября & \rule{0pt}{11mm} & \\
25 сентября & \rule{0pt}{11mm} & \\\bottomrule\end{tabularx}
\vspace{4mm}\begin{goalbox}\textbf{Решение sprint review:} $\square$ принят \quad $\square$ принят с долгом \quad $\square$ не принят.\\[2mm]\textbf{Незакрытые P0 и владелец:} \rule{105mm}{0.4pt}\end{goalbox}
\end{document}
